\section{The \texttt{ESS\_butterfly} McStas Component}
ESS butterfly moderator, 2016 revision

\subsection*{Identification}
\begin{itemize}
  \item \textbf{Site:} 
  \item \textbf{Author:} Peter Willendrup and Esben Klinkby
  \item \textbf{Origin:} DTU
  \item \textbf{Date:} August-September 2016
\end{itemize}

\subsection*{Description}
\begin{lstlisting}
ESS butterfly moderator with automatic choice of coordinate system, with origin
placed at relevant "Moderator Focus Coordinate System" depending on sector location.

To select beamport N 5 simply use

COMPONENT Source = ESS_butterfly(sector="N",beamline=5,Lmin=0.1,Lmax=20,dist=2,
cold_frac=0.5, yheight=0.03,focus_xw=0.1, focus_yh=0.1)

<b>Geometry</b>
The geometry corresponds correctly to the latest release of the butterfly moderator,
including changes warranted by the ESS CCB in July 2016, the so called BF1 type moderator.
A set of official release documents are available with this component, see the benchmarking
website mentioned below.

<b>Brilliances, geometry adapted from earlier BF2 design</b>
The geometry and brightness data implemented in the McStas ESS source component ESS_butterfly.comp,
are released as an updated component library for McStas 2.3, as well as a stand alone archive for
use with earlier versions of McStas.

The following features are worth highlighting:
<ul>
<li>The brightness data are still based on last years MCNP calculations, based on the Butterfly 2 geometry.
As a result, the spatial variation of the brightness across the moderator face should be considered to
have an uncertainty of the order of 10%. Detailed information on the reasoning behind the change to
the Butterfly 1 geometry can be found in <A HREF="http://essbutterfly.mcstas.org/PDFs/Update_to_ESS_moderators_KHA_latest.pdf">[1]</a> and detailed information on horizontal spatial brightness
variation can be found in <A HREF="http://essbutterfly.mcstas.org/PDFs/BFpaper_LZ_latest.pdf">[2]</a>. The spectral shape has been checked and has not changed significantly.
<li>A scaling factor has been introduced to in order to account for the decrease in brightness since 2015.
To accommodate the influence of the changed geometry, this scaling factor has been applied independently
for the cold and thermal contributions and is beamline dependent. It is adjusted to agree with the
spectrally-integrated 6cm width data shown in <A HREF="http://essbutterfly.mcstas.org/PDFs/Update_to_ESS_moderators_KHA_latest.pdf">[1]</a>,Figure 3.
<li>To allow future user adjustments of brilliance, the scalar parameters c_performance and t_performance
have been implemented. For now, we recommend to keep these at their default value of 1.0.
<li>The geometry has been updated to correspond within about 2 mm to the geometry described in <A HREF="http://essbutterfly.mcstas.org/PDFs/Update_to_ESS_moderators_KHA_latest.pdf">[1]</a>. This
has been done by ensuring that the position and apparent width of the moderators correspond to <A HREF="http://essbutterfly.mcstas.org/PDFs/Update_to_ESS_moderators_KHA_latest.pdf">[1]</a>,Figure 2,
which has been derived from current MCNP butterfly 1 model.
<li>The beamport is now defined directly by its sector and number (e.g. 'W' and '5'), rather than giving the angle,
as before. <A HREF="http://essbutterfly.mcstas.org/PDFs/Update_to_ESS_moderators_KHA_latest.pdf">[1]</a>,Figure 5 shows the geometry of the moderator2, beamport insert and beamline axis for beamline W5.
Since the underlying data is still from last years MCNP run, when the brightness was calculated at 10-degree
intervals, this means that the spectral curve for the nearest beamport on the grid 5,15,25,35,45,55 degrees
is used. The use of this grid has no effect on the accuracy of the geometry or brilliance because of the above-
mentioned beamline-dependent adjustments to the brilliance and geometry. See the website <A HREF="http://essbutterfly.mcstas.org/">[3]</a> for details.
</ul>
As before, the beamports all originate at the focal point of the sector. The beamline will in almost all cases be
horizontally tilted in order to view the cold or thermal moderator, which should be done using an Arm component.

<p>We expect to release an MCNP-event-based source model later in 2016, and possibly also new set of brilliance
functions for ESS_butterfly.comp. These are expected to include more realistic brilliances in terms of variation
across sectors and potentially also performance losses due to engineering reality. </b>

<b>Engineering reality</b>
An ad-hoc method for future implementation of "engineering reality" is included, use the
"c_performance/t_performance" parameters to down-scale performance uniformly across all wavelengths.

<b>References:</b>
<ol>
<li><A HREF="http://essbutterfly.mcstas.org/PDFs/Update_to_ESS_moderators_KHA_latest.pdf">Release document "Update to ESS Moderators, latest version"</a>
<li><A HREF="http://essbutterfly.mcstas.org/PDFs/BFpaper_LZ_latest.pdf">Release document "Description and performance of the new baseline ESS moderators, latest version"</a>
<li><A HREF="http://essbutterfly.mcstas.org/">http://essbutterfly.mcstas.org/</a> benchmarking website with comparative McStas-MCNP figures
<li><a href="http://essbutterfly.mcstas.org/visualisation">html-based, interactive 3D model of moderators and monolith, as seen from beamline N4</a>.
<li><A HREF="https://github.com/mccode-dev/McCode/blob/master/mcstas-comps/sources/ESS_butterfly.comp">Source code</A> for <CODE>ESS_butterfly.comp</CODE> at GitHub.
</ol>
\end{lstlisting}

\subsection*{Input parameters}
Parameters in \textbf{boldface} are required; the others are optional.

\begin{longtable}{p{0.22\textwidth}p{0.12\textwidth}p{0.46\textwidth}p{0.14\textwidth}}
\toprule
\textbf{Name} & \textbf{Unit} & \textbf{Description} & \textbf{Default} \\
\midrule
\endhead
sector & str & Defines the 'sector' of your instrument position. Valid values are "N","S","E" and "W" & "N" \\
beamline & 1 & Defines the 'beamline number' of your instrument position. Valid values are 1..10 or 1..11 depending on sector & 1 \\
yheight & m & Defines the moderator height. Valid values are 0.03 m and 0.06 m & 0.03 \\
cold\_frac & 1 & Defines the statistical fraction of events emitted from the cold part of the moderator & 0.5 \\
target\_index & 1 & Relative index of component to focus at, e.g. next is +1 this is used to compute 'dist' automatically. & 0 \\
dist & m & Distance from origin to focusing rectangle; at (0,0,dist) - alternatively use target\_index & 0 \\
focus\_xw & m & Width of focusing rectangle & 0 \\
focus\_yh & m & Height of focusing rectangle & 0 \\
c\_performance & 1 & Cold brilliance scalar performance multiplicator c\_performance \textgreater{} 0 & 1 \\
t\_performance & 1 & Thermal brilliance scalar performance multiplicator t\_performance \textgreater{} 0 & 1 \\
\textbf{Lmin} & AA & Minimum wavelength simulated &  \\
\textbf{Lmax} & AA & Maximum wavelength simulated &  \\
tmax\_multiplier & 1 & Defined maximum emission time at moderator, tmax= tmax\_multiplier * ESS\_PULSE\_DURATION. & 3 \\
n\_pulses & 1 & Number of pulses simulated. 0 and 1 creates one pulse. & 1 \\
acc\_power & MW & Accelerator power in MW & 5 \\
tfocus\_dist & m & Position of time focusing window along z axis & 0 \\
tfocus\_time & s & Time position of time focusing window & 0 \\
tfocus\_width & s & Time width of time focusing window & 0 \\
\bottomrule
\end{longtable}

\subsection*{Links}
\begin{itemize}
  \item \href{run:/home/willend/willend-McCode/mcstas-comps/sources/ESS_butterfly.comp}{Source code} for \texttt{ESS\_butterfly.comp}.
\end{itemize}
\IfFileExists{ESS_butterfly_static.tex}{\input{ESS_butterfly_static.tex}}{}